% Partie V condensée (FR). Trois tables de résultats + schéma du détour. En-têtes
% traduits, nombres en virgule décimale (comme la table de l'article CamemBERT-bio).

\noindent Nous entraînons ensuite des encodeurs biomédicaux compacts. Poursuivre le
pré-entraînement d'un modèle français généraliste sur \emph{biomed-fr} améliore les
cinq jeux de reconnaissance d'entités nommées biomédicales testés
(\Cref{tab:bioVSbase}), sans en entraîner un nouveau à partir de zéro, pour un gain
moyen de 2,54 points et pour une fraction du coût. CamemBERT-bio a été adapté pour
environ 0,80\,kg de CO\textsubscript{2}, contre 26,11 pour DrBERT entraîné à partir de
zéro~\citep{labrak2023drbert,touchent2023camembertbio}.

\begin{table}[ht]
\centering
\small
\setlength{\tabcolsep}{5pt}
\begin{threeparttable}[b]
\begin{tabular}{l c c c}
\toprule
& & \multicolumn{2}{c}{\textbf{CamemBERT-bio}} \\
\cmidrule(lr){3-4}
\textbf{Jeu} & \textbf{CamemBERT} & \textbf{biomed-fr-small} & \textbf{biomed-fr} \\
\midrule
\rowcolor{ThesisTableSep}
\multicolumn{4}{c}{\rule{0pt}{10pt}\textbf{Clinique}} \\[1pt]
CAS1 & 70,50 $\pm$ 1,75 & \underline{72,94 $\pm$ 1,12} & \textbf{73,03 $\pm$ 1,29} \\
CAS2 & 79,02 $\pm$ 0,92 & \underline{80,00 $\pm$ 0,32} & \textbf{81,66 $\pm$ 0,59} \\
E3C  & 67,63 $\pm$ 1,45 & \underline{67,96 $\pm$ 1,85} & \textbf{69,85 $\pm$ 1,58} \\
\midrule
\rowcolor{ThesisTableSep}
\multicolumn{4}{c}{\rule{0pt}{10pt}\textbf{Notices}} \\[1pt]
EMEA & 74,14 $\pm$ 1,95 & \underline{75,93 $\pm$ 2,42} & \textbf{76,71 $\pm$ 1,50} \\
\midrule
\rowcolor{ThesisTableSep}
\multicolumn{4}{c}{\rule{0pt}{10pt}\textbf{Scientifique}} \\[1pt]
MEDLINE & \underline{65,73 $\pm$ 0,40} & 65,48 $\pm$ 0,31 & \textbf{68,47 $\pm$ 0,54} \\
\midrule
\textbf{Moyenne} & 71,40 & \underline{72,46} & \textbf{73,94} \\
\bottomrule
\end{tabular}
\begin{tablenotes}
  \small
  \item F-mesure (micro-moyenne, seqeval strict, IOB2). Moyenne $\pm$ écart-type sur 10 amorces. biomed-fr-small = sous-ensemble de 10\%.
\end{tablenotes}
\end{threeparttable}
\caption{\textbf{CamemBERT vs CamemBERT-bio} sur cinq jeux de reconnaissance d'entités nommées biomédicales. Le pré-entraînement continu sur \textit{biomed-fr} apporte +2,54 de F-mesure en moyenne.}
\label{tab:bioVSbase}
\end{table}

Nous introduisons ensuite un détour causal : le modèle apprend par prédiction du
prochain token, puis revient au masquage bidirectionnel standard. Comme le montre la
\Cref{fig:clm-pipeline}, le détour laisse l'architecture inchangée et ne modifie que
le masque d'attention et la fonction de coût. Pendant la phase causale, le modèle ne
peut pas regarder en avant : il doit donc compresser l'information dans chaque
représentation de token, et la redescente ultérieure vers le masquage ne le ramène pas
à un état de type MLM.

\begin{figure}[ht]
\centering
\begin{tikzpicture}[node distance=1.8cm, >=Stealth, scale=0.85, transform shape,
  box/.style={rectangle, rounded corners=3pt, draw=ThesisNeutral, minimum height=0.9cm, minimum width=2.4cm, align=center, font=\small\sffamily}
]
  \node[box, fill=ThesisPrimary!20] (base) {encodeur MLM\\(ModernBERT)};
  \node[box, fill=ThesisTertiary!30, right=of base] (clm) {phase CLM\\(10--50\,Md tokens)};
  \node[box, fill=ThesisSecondary!30, right=of clm] (decay) {décroissance MLM\\(10\% de la Phase 1)};
  \node[box, fill=ThesisPrimary!30, right=of decay] (out) {ModernCamem-\\BERT-bio};

  \draw[->, thick, ThesisNeutral] (base) -- (clm) node[midway, above, font=\tiny\sffamily] {masque causal};
  \draw[->, thick, ThesisNeutral] (clm) -- (decay) node[midway, above, font=\tiny\sffamily] {masque bidir.};
  \draw[->, thick, ThesisNeutral] (decay) -- (out);
\end{tikzpicture}
\caption{\textbf{La chaîne du détour CLM.} À partir d'un encodeur MLM pré-entraîné, on poursuit le pré-entraînement avec un objectif causal (Phase~1), puis on redescend vers le MLM (Phase~2). Le modèle ne revient pas à un état de type MLM.}
\label{fig:clm-pipeline}
\end{figure}

Sur les huit tâches biomédicales françaises, le détour dépasse la référence MLM
appariée, comme le rapporte la \Cref{tab:french-results}.

\begin{table}[ht]
\centering
\small
\caption{\textbf{Résultats biomédicaux français} (F-mesure macro, 9 amorces). En gras : meilleur par colonne ; souligné : deuxième.}
\label{tab:french-results}
\resizebox{\textwidth}{!}{
\begin{tabular}{lcccccccccc}
\toprule
& & \multicolumn{5}{c}{\textbf{Classification multi-étiquettes}} & {\textbf{CLS}} & \multicolumn{2}{c}{\textbf{NER}} & \\
\cmidrule(lr){3-7} \cmidrule(lr){8-8} \cmidrule(lr){9-10}
& \textbf{Ctx} & \textbf{FRACCO-30} & \textbf{FRACCO-100} & \textbf{CANTEMIST} & \textbf{DISTEMIST} & \textbf{MedDialog} & \textbf{DiaMED} & \textbf{EMEA} & \textbf{Medline} & \textbf{Moy.} \\
\midrule
\rowcolor{ThesisTableSep} \multicolumn{11}{c}{\rule{0pt}{10pt}\textbf{Références}} \\[1pt]
ModernCamemBERT & 8192 & 70,1 & 55,3 & 63,3 & 20,2 & 60,6 & 56,4 & 68,0 & 59,7 & 56,7 \\
DrBERT & 512 & 53,0 & 35,6 & 37,9 & 21,4 & \underline{63,6} & 57,0 & 69,6 & 62,8 & 50,1 \\
CamemBERT-bio & 512 & 41,9 & 20,1 & 12,8 & 9,6 & 38,6 & 47,7 & \textbf{70,8} & \textbf{65,2} & 38,3 \\
CamemBERT & 512 & 40,8 & 19,4 & 11,9 & 9,5 & 37,4 & 40,6 & 69,5 & 62,7 & 36,5 \\
\midrule
\rowcolor{ThesisTableSep} \multicolumn{11}{c}{\rule{0pt}{10pt}\textbf{Nos modèles (Base, 150M)}} \\[1pt]
Référence MLM & 8192 & 69,9 & 56,8 & 64,9 & 23,5 & 62,5 & 63,4 & 68,5 & 61,4 & 58,9 \\
Détour CLM & 8192 & 74,8 & 60,1 & 71,0 & 25,5 & \underline{63,6} & \textbf{67,4} & 68,6 & 61,9 & 61,6 \\
\midrule
\rowcolor{ThesisTableSep} \multicolumn{11}{c}{\rule{0pt}{10pt}\textbf{Nos modèles (Large, 350M)}} \\[1pt]
Référence MLM & 8192 & \underline{79,4} & \underline{63,3} & \underline{72,6} & \underline{29,1} & \textbf{64,5} & 61,2 & \underline{70,4} & \underline{63,5} & \underline{63,0} \\
Détour CLM & 8192 & \textbf{80,7} & \textbf{65,4} & \textbf{74,4} & \textbf{30,4} & \textbf{64,5} & \underline{64,8} & 70,3 & 63,1 & \textbf{64,2} \\
\bottomrule
\end{tabular}
}
\end{table}

Il atteint 61,6\% de F-mesure moyenne contre 58,9\% en Base, et 64,2\% contre 63,0\%
en Large. Comme les deux trajectoires ne diffèrent que par l'objectif de la Phase~1,
l'écart est l'effet de cet objectif seul, sans coût supplémentaire, et il se transfère
à l'anglais. Nous publions ces modèles français sous le nom ModernCamemBERT-bio, MC-bio en abrégé.

Nous introduisons aussi \emph{OntoBook}, qui complète le pré-entraînement textuel en
transformant des ontologies médicales en manuels synthétiques : un sous-graphe
d'ontologie est converti en une marche aléatoire puis reformulé en prose fluide, et
l'encodeur s'entraîne dessus avec deux objectifs alignés, le masquage et la prédiction
de relation entre paires de codes. Il dépasse toutes les références sur les trois jeux
de codage français, avec le plus grand gain sur le codage fin, où il réduit
aussi la variance de $\pm$5,8 à $\pm$1,1.

Nous publions les modèles et les données obtenus.
